% 02_related_work.tex -- the rehauled paper's section, as bullets to expand into prose (2026-10-15).
% Sources and pointers behind every number: writeup/PAPER_REHAUL_OUTLINE_2026-10-15.md, section 2.
% House rules: every number in the final prose comes from a generated macro (writeup/generated/*.tex), never
% typed; park (comment out with a dated note) rather than delete.
% Macros to use here: cite only keys in references.bib; the [to add] items need bib entries
% 2026-10-15 (notation pass, writeup/NOTATION_GUIDE_2026-10-15.md): the literature's own terms; every \citet key
% checked against references.bib. Old -> new: \mathbb E_\tau RV -> \E_t \RV^{(k)} (daily RV averaged over k trading
% days, kept apart from the paper's closing-interval \RV_t); B_t -> s^2_{t,i}; h (forecast) -> v (h is the DM
% horizon); \beta -> \bm\beta; \lambda_1, \lambda_2 -> \lambda, \alpha (the guide's penalty and mixing; \lambda_1,
% \lambda_2 stay XGBoost's); g = pooled -> pooled intraday model.
% 2026-10-15 (finish, user): "Intraday periodicity" heading -> "Diurnal pattern"; "periodicity factor" -> "diurnal
% adjustment" (term rename only; s^2_{t,i} stays; Andersen-Bollerslev's own "intraday periodicity" kept); the
% retransformation (smearing) bullet parked.
% 2026-10-15 (finish): the Introduction is parked, so this section opens the paper; symbols used here are now defined
% at first use (E_t, the daily target RV^{(d)}_{t+1}, RV_{t,i} and s^2_{t,i}, the measures Q and P, RV_t, \bm x_t,
% QLIKE, \bm y, \bm X, lambda, phi_{t,j}); the MinGRU and PufferLib cites are parked (user: NN content parked).

% 2026-10-15 (consolidation): the HAR model in Corsi's own daily, weekly and monthly notation, so that the
% superscript (k) keeps one unit in the paper (k intraday intervals, the HAR horizon); X -> \bm X.
\section{Related work}
\label{sec:rh_related}
% 2026-10-17 (user): prose rewritten in the style of Zhang, Zhang, Cucuringu, and Qian (2024).
% 2026-10-24 (user, on E1): E1 parked text that exposited our estimators (components, target transform, origin);
% it does not mean Related work avoids our paper. Each bullet places our setup and results in its strand at the
% granularity an expert reader needs (what is the same, what is new, the obvious contrast), assuming the reader knows
% the literature; definitions of our estimators and trading rule stay in Methods and Economics. Framing pass below.
\begin{itemize}
% 2026-10-15 (consolidation): parked -- (k) in trading days here vs intraday intervals everywhere else.
% \item \textbf{Realized variance and the HAR model:}
%   $\E_t[\RV^{(1)}_{t+1}] = \beta_0 + \sum_{k} \beta_k\, \RV^{(k)}_t$, $k \in \{1,5,22\}$ trading days, with
%   $\RV^{(k)}_t$ the average daily realized variance over the last $k$ trading days. Our HAR components
%   $\bar y^{(k)}_t$ take the same averages of $y$ over $k$ intraday intervals.\newline
% 2026-10-15 (finish): parked -- the Introduction is parked, so Related work opens the paper; the bullet now defines
% E_t and the daily target and reads standalone. Was: "$\E_t[\RV_{t+1}] = \beta_0 + ...$, with $\RV^{(d)}_t$, ... the
% average daily realized variance over the last 1, 5 and 22 trading days. Our HAR components $\bar y^{(k)}_t$ take the
% same averages of $y$ over $k$ intraday intervals."
\item \textbf{Realized variance and the HAR model.} Following \citet{Corsi2009}, the heterogeneous
  autoregressive (HAR) model specifies the conditional expectation of next-day realized variance as a linear
  function of past realized variances aggregated over daily, weekly, and monthly horizons:
  $\E_t[\RV^{(d)}_{t+1}] = \beta_0 + \beta_d\, \RV^{(d)}_t + \beta_w\, \RV^{(w)}_t + \beta_m\, \RV^{(m)}_t$,
  where $\E_t$ denotes the conditional expectation given information available at the forecast origin on day $t$,
  and $\RV^{(d)}_t$, $\RV^{(w)}_t$, and $\RV^{(m)}_t$ denote the average daily realized variance over the
  preceding 1, 5, and 22 trading days, respectively. Because all three components represent daily realized
  variances aggregated at the daily, weekly, and monthly horizons, forecasting the realized variance of a
  specific intraday interval lies outside this standard daily specification. Our benchmark carries the HAR
  structure to the intraday interval, and the question we ask is how much a large set of exogenous predictors adds
  to it.\newline
  See \citet{Corsi2009}, \citet{AndersenBollerslevDieboldLabys2003}, \citet{BarndorffNielsenShephard2004},
  \citet{AndersenBollerslevDiebold2007}, \citet{CorsiMittnikPigorschPigorsch2008},
  \citet{BollerslevPattonQuaedvlieg2016}, \citet{BekiermanManner2018}, \citet{BennedsenLundePakkanen2022},
  \citet{HansenHuangShek2012}, and \citet{AndersenBollerslevChristoffersenDiebold2006}.
  % 2026-10-16 (user, E1): parked -- Related work does not describe our estimators. The components, the square-root
  % target, and the 15:30 origin belong in Methods. Was: "... on day $t$ (here 15:30 ET) and ... Our HAR components
  % $\bar y^{(k)}_t$ take the same averages over $k$ intraday intervals of $y$, the square root of realized variance
  % divided by its diurnal adjustment (Section~\ref{sec:rh_data})."
  % \editc{red}{E1}{Related work says the HAR components average $y$... Predictors and Methods average $\tilde y$.}
  \begin{itemize}
  % 2026-10-24 (user, framing pass): parked -- the sentence without the link to our pooling. Was: "In a multivariate
  % setting, the panel HAR framework of \citet{BollerslevHoodHussLebedev2018} estimates a single pooled regression
  % jointly across a cross-section of assets."
  \item In a multivariate setting, the panel HAR framework of \citet{BollerslevHoodHussLebedev2018} estimates a
    single pooled regression jointly across a cross-section of assets; our intraday models pool in the same spirit,
    across the intraday intervals of one index rather than across assets.
    % 2026-10-16 (user, E1): parked -- this named our estimator. Was: "the cross-asset counterpart of our intraday
    % model, which pools the intraday intervals of one index."
    % Key says "Lebedev"; the bib's author field (correctly) says Pedersen -- \citet prints the author field.
  \item For the theoretical foundations of realized variance, see \citet{BarndorffNielsenShephard2002}.
  % 2026-10-15 (bib pass): parked -- the [to add] placeholder, now a \citet key in references.bib. Was:
  % \item \todo{to add: Barndorff-Nielsen--Shephard 2002 (realized variance)}.
  \end{itemize}
\item \textbf{Diurnal pattern.} It is well documented that intraday return volatility exhibits a pronounced
  U-shaped pattern across the trading session. To account for this periodicity, standard practice scales intraday
  realized variance by a diurnal factor estimated from the corresponding clock-time interval on prior trading
  days. We follow this practice; it is what allows one model to be estimated across all intervals of the trading
  day.\newline
  % 2026-10-16 (user, E1): parked -- the display was our target, not the literature's. Was:
  % "divides by a diurnal adjustment, here $y_{t,i} = \sqrt{\RV_{t,i}/s^2_{t,i}}$, with $\RV_{t,i}$ the realized
  % variance of intraday interval $i$ on day $t$ and $s^2_{t,i}$ its diurnal adjustment."
  See \citet{AndersenBollerslev1997} on intraday periodicity, \citet{EngleGallo2006} on intraday range and
  volatility measures as indicators of daily volatility, and \citet{HizmeriBevilacqua2025} on morning
  volatility uncertainty.
  % 2026-10-15 (notation pass): EngleGallo2006 is a multiple-indicators model of daily volatility, not a periodicity
  % paper; its claim is narrowed to what it shows.
  \begin{itemize}
  \item See also \citet{TaylorXu1997} on diurnal volatility adjustments and \citet{EngleSokalska2012} on
    multiplicative decompositions of intraday volatility.
  % 2026-10-15 (bib pass): parked -- the [to add] placeholder, now \citet keys in references.bib. Was:
  % \item \todo{to add: Taylor--Xu 1997 (periodicity factor); Engle--Sokalska 2012 (intraday volatility)}.
  \end{itemize}
\item \textbf{0DTE options and the variance risk premium.} The variance risk premium, defined as the difference
  between option-implied variance and expected realized variance,
  $\E^{\mathbb{Q}}_t[\RV_t] - \E^{\mathbb{P}}_t[\RV_t]$, under the risk-neutral ($\mathbb Q$) and physical
  ($\mathbb P$) measures (where $\RV_t$ denotes the realized variance through expiration), is empirically
  positive on average. Our economic application sits at the end of a 0DTE option's life: we ask whether more
  accurate forecasts of the closing half hour's realized variance, set against the option-implied variance,
  translate into straddle returns.\newline
  See \citet{HizmeriBevilacqua2025}, as well as the foundational option-pricing frameworks of
  \citet{BlackScholes1973} and \citet{Merton1973}.
  \begin{itemize}
  \item On zero-days-to-expiration (0DTE) options, see \citet{BandiFusariReno2023},
    \citet{BeckmeyerBrangerGayda2023}, \citet{BrogaardHanWon2023}, \citet{DimErakerVilkov2024},
    \citet{Vilkov2026}, and \citet{AlmeidaFreireHizmeri2025}; on the variance risk premium more broadly, see
    \citet{CarrWu2009} and \citet{BollerslevTauchenZhou2009}.
  % 2026-10-15 (bib pass): parked -- the [to add] placeholder, now \citet keys in references.bib. Was:
  % \item \todo{to add: Bandi--Fusari--Ren\`o 2023; Beckmeyer--Branger--Gayda 2023; Brogaard--Han--Won 2023;
  %   Dim--Eraker--Vilkov 2024; Vilkov, 0DTE trading rules; Almeida--Freire--Hizmeri 2025; Carr--Wu 2009;
  %   Bollerslev--Tauchen--Zhou 2009}.
  \end{itemize}
\item \textbf{Machine learning for volatility and returns.} A growing literature models conditional return and
  volatility dynamics via flexible nonlinear mappings $f(\bm x_t)$ of a high-dimensional predictor vector
  $\bm x_t$. Most of this work targets daily volatility or monthly returns. The closest precedent for our
  setting, \citet{ZhangZhangCucuringuQian2024}, forecasts the realized volatility of intraday intervals for a
  cross-section of stocks with one model across the intervals of the day and predictors built from realized
  volatility alone; we keep that design for the S\&P~500 and ask what a large set of exogenous predictors adds,
  comparing penalized linear models with tree ensembles.\newline
  See \citet{gu2020empirical} for return prediction and \citet{christensen2023machine} for realized volatility
  forecasting; for gradient-boosted decision trees, see \citet{Ke2017} (LightGBM) and \citet{Chen2016}
  (XGBoost); for Bayesian hyperparameter optimization, see \citet{Akiba2019} (Optuna).
  \begin{itemize}
  % 2026-10-15 (user): parked -- NN content parked for now. Was: "...; realized volatility: \citet{Bucci2020},
  % \citet{RahimikiaPoon2020}, \citet{ZhangZhangCucuringuQian2024}; MinGRU \citet{Feng2024}; PufferLib \citet{Suarez2024}."
  \item For ensemble methods, see \citet{Breiman2001} on random forests and \citet{Friedman2001} on gradient
    boosting; for machine-learning applications to realized volatility and intraday commonality, see
    \citet{Bucci2020}, \citet{RahimikiaPoon2020}, and \citet{ZhangZhangCucuringuQian2024}.
  % 2026-10-15 (bib pass): parked -- the [to add] placeholder, now \citet keys in references.bib. Was:
  % \item \todo{to add: Breiman 2001; Friedman 2001 (gradient boosting); Bucci 2020; Rahimikia--Poon 2020;
  %   Zhang--Zhang--Cucuringu--Qian 2024; MinGRU (Feng et al.\ 2024); PufferLib}.
  \end{itemize}
% 2026-10-16 (user): parked -- forecast evaluation is standard. Cite Patton, Diebold--Mariano,
% Harvey--Leybourne--Newbold, Newey--West, Ehm, and the model-confidence set where those methods are used
% (Forecast evaluation, Statistical tests, Economic value). Mincer--Zarnowitz and White's reality check are not used.
% \item \textbf{Forecast evaluation:} QLIKE is robust to noise in the volatility proxy: with a conditionally unbiased
%   proxy it ranks forecasts as the true conditional variance would, and
%   $\argmin_{v > 0} \E_t[\QLIKE(\RV_t, v)] = \E_t[\RV_t]$, with
%   $\QLIKE(\RV, v) = \RV/v - \log(\RV/v) - 1$.\newline
%   \citet{Patton2011} (robust loss functions), \citet{DieboldMariano1995} with the small-sample correction of
%   \citet{HarveyLeybourneNewbold1997}, \citet{HansenLundeNason2011} (model confidence set), \citet{Ehm2016}
%   (Murphy diagrams, elementary scores), \citet{MincerZarnowitz1969} (forecast-efficiency regressions).
%   \begin{itemize}
%   \item \citet{Hansen2005} (superior predictive ability), \citet{White2000} (reality check), \citet{RomanoWolf2005}
%     (StepM), \citet{NeweyWest1987} (HAC variance).
%   % 2026-10-15 (bib pass): parked -- the [to add] placeholder, now \citet keys in references.bib. Was:
%   % \item \todo{to add: Hansen 2005 (superior predictive ability); White 2000 (reality check); Romano--Wolf 2005
%   %   (StepM); Newey--West 1987 (HAC variance)}.
%   \end{itemize}
\item \textbf{Shrinkage with many weak signals.} In high-dimensional linear regressions, penalized estimators
  regularize the coefficient vector via
  \[
    \hat{\bm\beta} = \argmin_{\bm\beta} \lVert \bm y - \bm X\bm\beta \rVert_2^2
      + \lambda\bigl(\alpha \lVert\bm\beta\rVert_1 + \tfrac{1-\alpha}{2} \lVert\bm\beta\rVert_2^2\bigr)
  \]
  for a response vector $\bm y$, a predictor matrix $\bm X$, and an overall penalty strength $\lambda$, nesting
  ridge regression ($\alpha = 0$) and the lasso ($\alpha = 1$) as special cases; multi-penalty ridge extensions
  assign distinct regularization weights across feature blocks via $\lambda \sum_j \omega_j \beta_j^2$, where the
  penalty weight $\omega_j$ is held constant within each block. Our predictor set is large and its signals
  individually weak, so the comparison of penalized linear models with tree ensembles, and of how their gains
  change with the length of the training window, bears directly on the dense-versus-sparse question.\newline
  See \citet{GiannoneLenzaPrimiceri2021} on dense versus sparse representations, \citet{ShenXiu2025} on
  learning from many weak signals, \citet{KozakNagelSantosh2020} on shrinking the cross-section of returns,
  \citet{HoerlKennard1970} on ridge regression, \citet{Tibshirani1996} on the lasso, \citet{ZouHastie2005} on the
  elastic net, and \citet{vandeWiel2021multipenalty} and \citet{vanWieringen2023ridge} on multi-penalty ridge
  estimation.
  \begin{itemize}
  \item See also \citet{Zou2006} on the adaptive lasso and feature-specific penalty weights, and
    \citet{FriedmanHastieTibshirani2010} on coordinate-descent estimation under the elastic-net
    parametrization.
  % 2026-10-15 (bib pass): parked -- the [to add] placeholder, now \citet keys in references.bib. Was:
  % \item \todo{to add: Zou 2006 (adaptive lasso, penalty weights); Friedman--Hastie--Tibshirani 2010 (elastic-net
  %   parametrization)}.
  \end{itemize}
% 2026-10-16 (user): parked -- SHAP is standard. The decomposition and the Lundberg cites stay in Variable
% importance, where the values are computed. Aas et al. (dependent features) is not used.
% \item \textbf{Variable importance:} Shapley additive explanations (SHAP) decompose each forecast additively,
%   $\hat y_t = \phi_0 + \sum_{j=1}^{p} \phi_{t,j}$, with $\phi_{t,j}$ the contribution of predictor $j$ of $p$ and
%   $\phi_0$ the average forecast.\newline
%   \citet{lundberg2017unified} (SHAP), \citet{lundberg2018consistent} (TreeSHAP), \citet{aas2021explaining}
%   (dependent features).
% 2026-10-15 (user): parked -- no smearing mentions
% \item \textbf{Retransformation} (one sentence): a smearing estimate would correct $\hat y_t^{\,2} s^2_{t,c}$ for
%   retransformation bias; we do not.\newline
%   \citet{Duan1983}, \citet{ManningMullahy2001}.
\end{itemize}

